

%%%%%%%%%%%%%%%%%%%%%%%%%%%%%%%%%%%%%%%%%%%%%%%%%%%%%%%
\section{Introduction à l'optimisation des performances}
%%%%%%%%%%%%%%%%%%%%%%%%%%%%%%%%%%%%%%%%%%%%%%%%%%%%%%%
\begin{tsp}{\'Etapes du traitement d'une requête}{

Toute requête SQL est traitée en trois étapes\,:

\begin{enumerate}
  \item \textbf{Analyse et traduction} de la requête. On vérifie
            qu'elle est correcte, et on l'exprime sous forme
             d'opérations.

  \item \textbf{Optimisation}\,: comment agencer au mieux les opérations,
              et quels algorithmes utiliser. On obtient un \textbf{plan
              d'exécution}.

    \item \textbf{Exécution de la requête}\,: le plan d'exécution est compilé
            et exécuté.
\end{enumerate}
}\end{tsp}
%%%%%%%%%%%%%%%%%%%%%%%%%%%%%%%%%%%%%%%%%%%%%%%%%%%%%%%
\begin{tsp}{Support du traitement d'une requête}{

Le traitement s'appuie sur les éléments suivants\,:

\begin{enumerate}
  \item \textbf{Le schéma de la base}, description
          des tables et chemins d'accès (dans le catalogue)

  \item \textbf{Des statistiques}\,: taille des tables, des index,
              distribution des valeurs

    \item \textbf{Des algorithmes}\,: il peuvent
               différer selon  les systèmes
\end{enumerate}

\textbf{Important}\,: on suppose que le temps d'accès à
ces informations est négligeable par rapport à l'exécution
de la requête.
}\end{tsp}
%%%%%%%%%%%%%%%%%%%%%%%%%%%%%%%%%%%%%%%%%%%%%%%%%%%%%%%
\begin{tsp}{L'optimisation sur un exemple}
{
Consid\'erons le sch\'ema\,:

\begin{itemize}
  \item $CINEMA (Cin\acute{e}ma, Adresse, G\acute{e}rant)$\\
   \item $SALLE (Cin\acute{e}ma, NoSalle, Capacit\acute{e})$
\end{itemize}
avec les hypothèses\,:
\begin{enumerate}
\item Il y a 300 n-uplets dans CINEMA, occupant 30 pages.
\item Il y a 1200 n-uplets dans SALLE, occupant 120 pages.
\end{enumerate}
}
\end{tsp}

\begin{tsp}{Expression d'une requ\^ete}
{
On considère  la requ\^ete\,:
{\it Adresse des cinémas ayant des salles
de plus  de 150 places}

En SQL, cette requ\^ete s'exprime de la manière suivante\,:\\

{\tt
\begin{tabbing}
\=SELECT \=Adresse\\
\>FROM   \>CINEMA, SALLE\\
\>WHERE  \>capacité $>$ 150\\
\> AND   \> CINEMA.cinéma = Salle.cinéma
\end{tabbing}
}
}
\end{tsp}

\begin{tsp}{En algèbre relationnelle}
{
Traduit en algèbre, on a plusieurs possibilit\'es. En voici deux\,:
\begin{enumerate}
\item    $ \pi_{Cinema} (\sigma_{Capacite > 150}
        (CINEMA \Join SALLE))$
\item $ \pi_{Cinema}(CINEMA \Join
    \sigma_{Capacite>150}(SALLE))$
\end{enumerate}

Soit une jointure suivie d'une s\'election, ou l'inverse.\\
\vspace*{0.4cm}

NB\,: on peut les représenter comme des arbres.
}
\end{tsp}

\begin{frame}
\frametitle{\'Evaluation des co\^uts (très simplifiée)}

On suppose qu'il n'y a que 5\,\% de salles de plus de 150
places.

\begin{enumerate}
\item Jointure\,: on
lit 3\,600 pages (120x30); Sélection\,: on obtient 5\,\% de 120 pages, 
soit 6 pages.\\ 

Nombre d'E/S\,: 3\,600 + 120x2 + 6 = 3\,846.

\item S\'election\,: on lit 120 pages et on
obtient 6 pages. Jointure\,:
on lit 180 pages (6x30) et on obtient 6 pages.\\
Nombre d'E/S\,: 120 + 6 + 180 + 6 = 312.
\end{enumerate}

$\Rightarrow$ la deuxième strat\'egie est de loin la meilleure\,!
\end{frame}

\begin{comment}
\begin{tsp}{Optimisation\,: principes généraux}
{
\begin{enumerate}
\item Il faut {\bf traduire} une requ\^ete exprim\'ee avec un langage
d\'eclaratif en une suite d'op\'erations (typiquement
les op\'erateurs de l'algèbre relationnelle).

\item En fonction (i) des co\^uts de chaque op\'eration (ii)
des caract\'eristiques de la base, (iii) des algorithmes
utilis\'es, on cherche \`a estimer la meilleure strat\'egie.

\item On obtient
le {\bf plan d'ex\'ecution} de la requ\^ete. Il
n'y a plus qu'\`a le traiter au niveau physique.
\end{enumerate}
}
\end{tsp}


\begin{tsp}{Les paramètres de l'optimisation}
{
Comme on l'a vu sur l'exemple, l'optimisation s'appuie sur\,:
\begin{enumerate}
\item Des {\bf règles de r\'e\'ecriture} des expressions de l'algèbre.
\item Des connaissances sur {\bf l'organisation physique} de la base 
(index, hachage,\ldots)
\item Des {\bf statistiques} sur les caract\'eristiques de
la base (taille des relations par exemple).
\end{enumerate}

Un {\bf modèle de co\^ut} permet de classer les diff\'erentes
strat\'egies envisag\'ees 
}
\end{tsp}

\begin{tsp}{Traduction alg\'ebrique}
{
D\'eterminer l'expression alg\'ebrique
\'equivalente \`a la requ\^ete\,:

\begin{enumerate}
\item  arguments du SELECT\,:  projections.
\item arguments du WHERE\,: $NomAttr1=NomAttr2$ correspond en g\'en\'eral
\`a une jointure, $NomAttr=constante$ \`a une s\'election.
\end{enumerate}

On obtient une expression alg\'ebrique qui peut \^etre repr\'esent\'ee
par un {\bf arbre de requ\^ete}. 
}
\end{tsp}

\begin{tsp}{Traduction alg\'ebrique\,: exemple}
{
Consid\'erons l'exemple suivant\,:\\
\centerline{\textit{Quels films passent au
REX à 20 heures\,?}}

{\tt\small
\begin{tabbing}
\=SELECT  \=film\\
\>FROM    \>Cinéma, Salle, Séance\\
\>WHERE   \>Cinéma.nom-cinéma = 'Le Rex'\\
\>AND     \>Séance.heure-début = 20\\
\>AND     \>Cinéma.nom-cinéma = Salle.nom-cinéma\\
\>AND     \>Salle.salle = Séance.salle
\end{tabbing}
}

$\pi_{film} (\sigma_{Nom='Le\,REX' \wedge heure-d\acute{e}but=20}
        ((CINEMA \Join SALLE) \Join SEANCE))$
}
\end{tsp}


\begin{tsp}{Restructuration}
{
Il y a plusieurs expressions
{\bf \'equivalentes} pour une m\^eme requ\^ete.\\


ROLE  DE L'OPTIMISEUR
\begin{enumerate}
  \item Trouver les expressions \'equivalentes \`a une requ\^ete.
  \item Les \'evaluer et choisir la meilleure.
\end{enumerate}

On convertit une expression en une expression \'equivalente
en employant des {\bf règles de réécriture}.
}
\end{tsp}

\begin{tsp}{Exemples de règles de réécriture}
{
\begin{enumerate}
\item {\bf Commutativit\'e des jointures}\,:\\
        $ R \Join S  \equiv S \Join R$
\item {\bf Associativit\'e des jointures}\,:\\
        $(R \Join S) \Join T \equiv R \Join (S \Join T)$
\item {\bf Regroupement des s\'elections}\,:\\
        $\sigma_{A='a' \wedge B='b'}(R) \equiv
        \sigma_{A='a'}(\sigma_{B='b'}(R))$
\item {\bf Commutativit\'e de la s\'election et de la projection}\\
%               $\pi_{A_1, A_2, \ldots A_p}(\sigma_{A_i='a} (R)) \equiv
%               \pi_{A_1, A_2, \ldots A_p}(\sigma_{A_i='a}( 
%               \pi_{A_1, A_2, \ldots A_p, A_i}(R)))$
        $\pi_{A_1, A_2, \ldots A_p}(\sigma_{A_i='a} (R)) \equiv
        \sigma_{A_i='a}(\pi_{A_1, A_2, \ldots A_p}(R)), i 
        \in \{1,\ldots,p\}$
\end{enumerate}
}
\end{tsp}
\begin{tsp}{Règles de réécriture (suite)}
{
\begin{enumerate}
\item {\bf Commutativit\'e de la s\'election et de la jointure}.\\
$\sigma_{A='a'} (R(\ldots A\ldots) \Join S) \equiv
     \sigma_{A='a'}(R) \Join S$
\vfill
%\item {\bf Distributivit\'e de la s\'election sur l'union}.\\
% $\sigma_{A='a'} (R \cup S) \equiv 
%     \sigma_{A='a'} (R) \cup \sigma_{A='a'} (S)$\\
%NB\,: valable aussi pour la diff\'erence.
\item  {\bf Commutativit\'e de la projection et de la jointure}\\
 $\pi_{A_1 \ldots A_pB_1\ldots B_q}(R \Join_{A_i=B_j} S) \equiv$\\
        $\pi_{A_1 \ldots A_p}(R) \Join_{A_i=B_j}\pi_{B_1\ldots B_q}(S))
  \;i \in \{1,\ldots ,p\}, j \in \{1,\ldots,q\}$
%\item {\bf Distributivit\'e de la projection sur l'union}
% $\pi_{A_1A_2\ldots A_p} (R \cup S) \equiv 
%     \pi_{A_1A_2\ldots A_p} (R) \cup \pi_{A_1A_2\ldots A_p} (S)$
\end{enumerate}
}
\end{tsp}


\begin{tsp}{Exemple d'un algorithme}
{
\begin{enumerate}
\item S\'eparer les s\'elections avec plusieurs pr\'edicats
en plusieurs s\'elections \`a un pr\'edicat (règle 3).
\item Descendre les s\'elections le plus bas possible dans l'arbre 
(règles 4, 5, 6).
\item Regrouper les s\'elections sur une m\^eme relation 
(règle 3).
\item Descendre les projections le plus bas possible (règles
7 et 8).
\item Regrouper les projections sur une m\^eme relation.
\end{enumerate}
}
\end{tsp}
%%%%%%%%%%%%%%%%%%%%%%%%%%%%%%%%%%%%%%%%%%%%%%%%%%%%%%%
\begin{tsp}{Que cherche-t-on à optimiser\,?}
{
On peut en fait définir plusieurs critères\,:

\begin{enumerate}
  \item \textbf{Temps d'exécution}\,: temps total d'exécution,
           souvent utilisé comme critère par défaut

  \item \textbf{Temps de réponse}\,: temps pour obtenir la
           première ligne du résultat. \`A considérer pour les
                     traitements en tâche de fond.

    \item \textbf{Débit transactionnel}\,: recherche de la fluidité
                 des traitements concurrents.

    \item \textbf{Transfert réseau}\,: pour les applications distribuées
\end{enumerate}

}\end{tsp}
\end{comment}
%%%%%%%%%%%%%%%%%%%%%%%%%%%%%%%%%%%%%%%%%%%%%%%%%%%%%%%
\begin{tsp}{En résumé}
{
Un module du SGBD, \textbf{l'optimiseur}, est chargé de\,:

\begin{enumerate}
  \item Prendre en entrée une requête, et la mettre sous forme d'opérations

  \item Se fixer comme objectif l'optimisation d'un certain 
           paramètre (en général le temps d'exécution)

    \item construire un programme s'appuyant sur
         les index existants, et les opérations disponibles.
\end{enumerate}

}\end{tsp}
%%%%%%%%%%%%%%%%%%%%%%%%%%%%%%%%%%%%%%%%%%%%%%%%%%%%%%%
\begin{tsp}{Le schéma de la base}{
\begin{itemize}
  \item Film (\cle{idFilm}, titre, année, genre, résumé,
                                \fk{idMES}, \fk{codePays})
 \item Artiste (\cle{idArtiste}, nom, prénom, annéeNaissance)
  \item Role (\cle{\fk{idActeur}, \fk{idFilm}}, nomRôle)
   \item Internaute (\cle{email}, nom, prénom, région)
  \item Notation (\cle{\fk{email}, \fk{idFilm}}, note)
   \item Pays (\cle{code}, nom, langue)
\end{itemize}

}\end{tsp}

\section{Plans d'exécution}
%%%%%%%%%%%%%%%%%%%%%%%%%%%%%%%%%%%%%%%%%%%%%%%%%%%%%%%
\begin{tsp}{L'essentiel de ce qu'il faut savoir}
{
Qu'est-ce qu'un plan d'exécution\,?
  \begin{itemize}
      \item C'est un \textbf{programme} combinant 
               des opérateurs physiques (chemins d'accès
                     et traitements de données).

      \item Il a la forme d'un \textbf{arbre}\,: chaque
              n{\oe}ud est un opérateur qui
  \begin{itemize}
      \item prend des données en entrée
      \item applique un traitement
      \item produit les données traitées en sortie
   \end{itemize}
\end{itemize}


}\end{tsp}
%%%%%%%%%%%%%%%%%%%%%%%%%%%%%%%%%%%%%%%%%%%%%%%%%%%%%%%
\begin{tsp}{L'essentiel de ce qu'il faut savoir (suite)}
{
La phase d'optimisation proprement dite\,:

  \begin{itemize}
     \item Pour \textbf{une} requête, le système a le choix
               entre \textbf{plusieurs} plans d'exécution.

     \item Ils diffèrent par l'ordre des opérations, 
             les algorithmes, les chemins d'accès.

     \item Pour chaque plan on peut estimer\,:
      \begin{itemize}
          \item le coût de chaque opération
         \item la taille du résultat
      \end{itemize}
      
\end{itemize}

\textbf{Objectif}\,: diminuer le plus vite possible la taille
des données manipulées.

}\end{tsp}
%%%%%%%%%%%%%%%%%%%%%%%%%%%%%%%%%%%%%%%%%%%%%%%%%%%%%%%
\begin{tsp}{Laissons le choix au système\,!}
{
Bon à savoir\,: il y a autant de plans d'exécution que
de \alert{blocs} dans une requête.


Exemple\,: cherchons tous les films avec James Stewart, parus
en 1958.\\

  \fichier{QStewart1.sql}

Pas d'imbrication\,: un bloc, OK\,!
}\end{tsp}
%%%%%%%%%%%%%%%%%%%%%%%%%%%%%%%%%%%%%%%%%%%%%%%%%%%%%%%
\begin{tsp}{Seconde requête (2 blocs)}{

La même, mais avec un niveau d'imbrication.

  \fichier{QStewart2.sql}

Une imbrication sans nécessité\,: moins bon\,!
}\end{tsp}
%%%%%%%%%%%%%%%%%%%%%%%%%%%%%%%%%%%%%%%%%%%%%%%%%%%%%%%
\begin{tsp}{Troisième requête (2 blocs)}{

La même, mais avec \texttt{EXISTS} au lieu de \texttt{IN}.\\

  \fichier{QStewart3.sql}

}\end{tsp}
%%%%%%%%%%%%%%%%%%%%%%%%%%%%%%%%%%%%%%%%%%%%%%%%%%%%%%%
\begin{tsp}{Quatrième requête (3 blocs)}{

La même, mais avec deux imbrications\,:\\

  \fichier{QStewart4.sql}

Très mauvais\,: on force le plan d'exécution, et il
est très inefficace.

}\end{tsp}
%%%%%%%%%%%%%%%%%%%%%%%%%%%%%%%%%%%%%%%%%%%%%%%%%%%%%%%
\begin{tsp}{Pourquoi c'est mauvais}{

\begin{itemize}
  \item On parcourt tous les films parus en 1958

  \item Pour chaque film\,: on cherche les rôles
       du film, \textbf{mais pas d'index disponible}
 
  \item Ensuite, pour chaque rôle on regarde si c'est James Stewart
\end{itemize}

Ca va coûter cher\,!!
}\end{tsp}
%%%%%%%%%%%%%%%%%%%%%%%%%%%%%%%%%%%%%%%%%%%%%%%%%%%%%%%
\begin{tsp}{Exemples de plans d'exécution}{

Gardons la même requête. Voici les opérations disponibles\,:\\

\figWithSize{algphys}{9}

}\end{tsp}
%%%%%%%%%%%%%%%%%%%%%%%%%%%%%%%%%%%%%%%%%%%%%%%%%%%%%%%
\begin{tsp}{Sans index sur le nom}{
\figWithSize{PE1}{10}

}\end{tsp}
%%%%%%%%%%%%%%%%%%%%%%%%%%%%%%%%%%%%%%%%%%%%%%%%%%%%%%%
\begin{tsp}{Avec index sur le nom}{
\figWithSize{PE2}{10}

}\end{tsp}
%%%%%%%%%%%%%%%%%%%%%%%%%%%%%%%%%%%%%%%%%%%%%%%%%%%%%%%
\begin{tsp}{Sans index}{
\figWithSize{PE4}{9}

}\end{tsp}

